
\chapter{Diseño experimental}
\label{sec:experimento}

Este capítulo detalla cómo se han diseñado y ejecutado los experimentos que
responden a RQ2 y RQ3. En primer lugar se describe el corpus sobre el que se
trabaja. A continuación se definen las métricas y la línea base con las que se
mide el efecto de la transformación. Después se presentan, por separado, el
protocolo determinista de RQ2 y el protocolo de evaluación de los modelos de
lenguaje de RQ3, y el capítulo se cierra con la validación complementaria de la
métrica frente a SonarQube. Un principio recorre todo el capítulo: la
\textbf{reproducibilidad}. Cada decisión experimental se fija de antemano y se
registra, de modo que el procedimiento aquí descrito pueda repetirse y conducir
al mismo resultado.

\section{Corpus}
\label{sec:corpus}
El corpus se carga, por defecto, desde recursos empaquetados mediante
\texttt{PilotCorpusLoader} y \texttt{RealDatasetLoader}; ambos cargadores
admiten además un \emph{directorio del sistema de archivos} ---con subcarpetas
\texttt{pilot-corpus/} y \texttt{real-corpus/}---, de modo que el corpus puede
sustituirse, sin recompilar, indicando esa ruta al ejecutar el lote. Su
composición se resume en la \autoref{tab:corpus}.

\begin{table}[ht]
\centering
\caption{Composición del corpus de los experimentos RQ1/RQ2.}
\label{tab:corpus}
\begin{tabular}{@{}lrl@{}}
\toprule
\textbf{Subconjunto} & \textbf{Casos} & \textbf{Finalidad} \\
\midrule
Piloto (sintético) & 8   & Casos canónicos y un rechazo por cada precondición \\
Real (código abierto) & 118 & Validación sobre código real de proyectos de código abierto \\
\midrule
\textbf{Total} & \textbf{126} & 111 elegibles, 15 no elegibles \\
\bottomrule
\end{tabular}
\end{table}

\paragraph{Corpus piloto}:

Ficheros sintéticos que cubren los casos canónicos del patrón objetivo y
sus rechazos típicos (casos válidos, inválidos por violación de cada
precondición y un caso mixto con varias oportunidades). Su finalidad es
servir de prueba controlada del pipeline.

\paragraph{Corpus real}:

El corpus real está formado por fragmentos extraídos de proyectos Java de código abierto. Cada caso conserva metadatos que permiten trazar su procedencia y contexto original. No son anotaciones de Java ni proceden de ninguna biblioteca: son etiquetas propias del proyecto, escritas como comentarios de línea al inicio de cada fichero (\lstinline|@project|, \lstinline|@file|, \lstinline|@method|, \lstinline|@license| y \lstinline|@sonarCCBefore|), que el cargador del corpus lee y asocia a cada caso.

Los proyectos fuente proceden de dos orígenes principales:

\begin{enumerate}
    \item \textbf{Proyectos de código abierto de uso general}, añadidos en este trabajo para ampliar la
    variedad de dominios y estilos de código. Aportaron casos Apache Ant,
    Apache POI, Jackson Core y Jackson Databind (FasterXML) y las bibliotecas
    Apache Commons (\emph{Lang}, \emph{Collections}, \emph{Math} e \emph{IO}).
    Otros proyectos de este grupo se escanearon pero no aportaron ningún caso
    (Apache Commons \emph{Numbers}, \emph{Geometry} y \emph{Compress}, ASM y
    PDFBox). El escaneo se realizó en varias tandas.

    \item \textbf{Proyectos de la evaluación de Saborido et
    al.}~\cite{10.1109/ACCESS.2022.3144743}: los diez proyectos del antecedente.
    Aportaron casos \texttt{cybercaptor-server}, \texttt{MOEAFramework},
    \texttt{knowage-server}, \texttt{fastjson}, \texttt{iotbroker}, jMetal y
    Bytecode Viewer; en cambio, \texttt{fiware-commons}, \texttt{jedis} y
    \texttt{aioteslake} se escanearon pero no aportaron ninguno.
\end{enumerate}

El detalle por proyecto (dominio, volumen de código analizado y casos
elegibles y no elegibles) se presenta junto al análisis de RQ1
(\autoref{sec:resultados-rq1}).

El corpus real no procede de un
muestreo arbitrario, sino del escaneo sistemático de los proyectos citados. En
conjunto, el escaneo evaluó 82\,655 candidatos a \texttt{if} externo del patrón
en 23 proyectos de código abierto y localizó 117 instancias elegibles del patrón
\texttt{if}--\texttt{if} combinable. A partir de esas instancias (junto con
casos no elegibles representativos y los 8 casos del piloto sintético) se
consolidó el corpus de 126 casos, de los cuales 99 son casos reales elegibles.
El corpus es, por tanto, esencialmente un censo de las oportunidades del patrón
presentes en los proyectos analizados, y no una submuestra de los
\emph{~}1\,000 métodos con $\text{CC}>15$ considerados en el anteproyecto.

\paragraph{Nota de alcance respecto al anteproyecto.} El anteproyecto aprobado
preveía partir de los \emph{~}1\,000 métodos con $\text{CC}>15$ del repositorio
de Saborido et al., sin acotar el tipo de transformación. Al concretar el
alcance en la combinación de condicionales anidados, un patrón específico y
verificable, no la reducción de CC en general, el corpus resultante quedó en
126 casos: los métodos que contienen exactamente este patrón en los proyectos
fuente, más los 8 casos sintéticos del piloto. El resto de los métodos con
$\text{CC}>15$ de esos proyectos pertenece a familias de transformación
distintas (principalmente extracción de método) que quedan fuera del alcance
declarado en \autoref{sec:alcance}.


\paragraph{Cobertura de RQ3.} La evaluación de RQ3 se ejecuta sobre el
\textbf{corpus consolidado completo} (126 casos: 8 piloto y 118 reales; 104
elegibles y 22 no elegibles), de modo que cubre todos los ejes de variación
(casos elegibles, parciales y no elegibles) y los principales motivos de descarte estructurales.
La \autoref{tab:subset-rq3} recoge una selección representativa de casos.

\begin{table}[ht]
\centering
\footnotesize
\caption{Ejemplos representativos de los casos de evaluación de RQ3
(la campaña cubre el corpus completo de 126 casos).}
\label{tab:subset-rq3}
\begin{tabularx}{\linewidth}{@{}lll>{\raggedright\arraybackslash}X@{}}
\toprule
\textbf{Identificador} & \textbf{Origen} & \textbf{Categoría} & \textbf{Justificación} \\
\midrule
\texttt{PILOT\_VALID\_SIMPLE} & Piloto & Elegible & Caso base \lstinline|if + if|. \\
\texttt{PILOT\_VALID\_NESTED\_IN\_LOOP} & Piloto & Elegible & Anidamiento en \lstinline|for|. \\
\texttt{REAL\_COMMONS\_MATH\_CONVERGED} & Real & Elegible & Doble condición numérica. \\
\texttt{REAL\_ANT\_MATCH\_PATH} & Real & Elegible & Triple \lstinline|if|; se combinan los dos pares anidados. \\
\texttt{REAL\_COMMONS\_MATH\_VALIDATE\_RANGE} & Real & No elegible (P1) & \lstinline|outer if| con \lstinline|else|. \\
\texttt{REAL\_COMMONS\_COLLECTIONS\_GET} & Real & Elegible & Llamada a método en la condición (combinable). \\
\bottomrule
\end{tabularx}
\end{table}

\section{Métricas y línea base}
La evaluación se apoya en dos métricas de complejidad y una referencia de
comparación, todas fijadas de antemano para garantizar la reproducibilidad:
\begin{itemize}
  \item \textbf{Métrica primaria:} complejidad cognitiva proxy (CC) calculada
  por \texttt{Cognitive\allowbreak Complexity\allowbreak Calculator} antes y después de la
  transformación; se reporta \emph{before}, \emph{after} y
  $\Delta = \text{after} - \text{before}$ (un $\Delta$ negativo indica
  reducción). Las limitaciones del proxy están en \autoref{sec:proxy}.
  \item \textbf{Métrica complementaria:} complejidad cognitiva medida por
  SonarQube Community Build sobre el corpus completo, que sirve para validar el
  proxy (\autoref{sec:sonar-val}).
  \item \textbf{Baseline determinista:} para cada caso, el resultado del
  pipeline (\texttt{BatchRunner}) constituye el baseline contra el que se
  contrastan las salidas de los LLMs en RQ3; declara explícitamente la
  elegibilidad y, en su caso, el $\Delta$ aplicable.
\end{itemize}

\section{Validación del proxy frente a SonarQube}
\label{sec:sonar-val}
Como la métrica primaria es un \emph{proxy} y no la implementación oficial,
conviene comprobar su validez antes de emplearla en los experimentos. Para ello
se contrastan sus valores con los de SonarQube Community Build (versión
\texttt{26.4.0.121862}) sobre el \textbf{corpus completo}.

El contraste se realiza caso por caso: se materializan los métodos \emph{before}
y \emph{after} como ficheros Java independientes, se analizan en dos proyectos
de SonarQube y se recupera su complejidad cognitiva oficial por fichero mediante
la API, que se compara con la del proxy.

El procedimiento es reproducible desde PowerShell: el análisis se lanza con el
\emph{goal} \texttt{mvn sonar:sonar} contra un servidor local de SonarQube,
arrancado en un contenedor \emph{Docker} mediante \texttt{docker-compose}. Los
resultados de este contraste se presentan en la \autoref{sec:resultados-sonar}.

\section{Protocolo RQ2 — Batch determinista \emph{before/after}}
\label{sec:exp-rq2}
El protocolo consta de cuatro pasos:

\begin{enumerate}
    \item Cargar la totalidad del corpus piloto y real.
    \item Para cada \texttt{ExperimentCase}, ejecutar
    \lstinline|BatchRunner.processCase| con el detector en modo
    \texttt{STRUCTURAL} (precondiciones estructurales P1--P4) y la política
    una por una.
    \item Recoger métricas \emph{before}, \emph{after} y $\Delta$, y motivos de
    descarte para los no elegibles.
    \item Exportar a JSON y CSV.
\end{enumerate}

De este protocolo se \textbf{congelan} ---es decir, se fijan de antemano y no
varían entre ejecuciones--- cuatro elementos: la composición del corpus, el modo
de detección (\texttt{STRUCTURAL}), el conjunto de precondiciones P1--P4 y la
política de aplicación una por una. Cualquier cambio en
\texttt{Cognitive\allowbreak Complexity\allowbreak Calculator} invalidaría las métricas y exigiría
repetir la ejecución completa.

\section{Protocolo RQ3 — Evaluación de LLMs}
\label{sec:exp-rq3}
El protocolo está congelado en
\lstinline|LlmExperimentProtocol.defaultProtocol()|
(\autoref{tab:protocolo-rq3}).

\begin{table}[ht]
\centering
\caption{Protocolo congelado de RQ3.}
\label{tab:protocolo-rq3}
\begin{tabular}{@{}lll@{}}
\toprule
\textbf{Parámetro} & \textbf{Valor} & \textbf{Justificación} \\
\midrule
Modelos & \texttt{gpt-4o}, \texttt{gpt-4.1} & Comparativa intra-proveedor (OpenAI). \\
Temperatura & \texttt{0.0} & Determinismo esperado. \\
Intentos por caso & \texttt{3} & Detección de inconsistencia residual. \\
Prompt & \texttt{v1.0} & Trazabilidad (texto literal en código). \\
\emph{Max tokens} salida & \texttt{2048} & Métodos refactorizados completos. \\
Cobertura & 126 casos & Corpus completo (elegibles/parciales/no elegibles). \\
Total invocaciones & 756 & $126 \times 2 \times 3$. \\
\bottomrule
\end{tabular}
\end{table}

La comparación se restringe deliberadamente a un único proveedor (OpenAI) por
tres razones. Primera, de \emph{validez interna}: contrastar dos modelos del
mismo proveedor (\texttt{gpt-4o} y \texttt{gpt-4.1}) mantiene constantes la
API, el tokenizador y la vía de acceso, de modo que las diferencias observadas
se atribuyen al modelo y no al entorno. Segunda, de \emph{alcance y recursos}:
el trabajo dispone de una única clave de API financiada y de un presupuesto
temporal acotado, lo que hace inviable una campaña multiproveedor sin
comprometer la profundidad del resto del estudio. Tercera, de
\emph{representatividad}: los modelos elegidos pertenecen a la familia de
propósito general de referencia en generación de código (linaje de Codex,
\autoref{sec:llm-antecedentes}). La generalización a otras familias de modelos
queda explícitamente como amenaza a la validez externa y trabajo futuro
(\autoref{cap:amenazas}); el diseño lo facilita, pues el \emph{provider} del
protocolo es intercambiable y añadir otro proveedor es un cambio de
configuración, no de arquitectura.

El procedimiento que se aplica a cada caso consta de los siguientes pasos:
\begin{enumerate}
  \item cargar el caso y su \emph{baseline} determinista;
  \item construir el \emph{prompt} \texttt{v1.0};
  \item invocar el modelo;
  \item almacenar la respuesta cruda;
  \item aplicar el oráculo (\autoref{fig:oraculo});
  \item registrar el \texttt{LlmEvaluationResult}.
\end{enumerate}
Estos pasos se repiten para cada intento y cada modelo. Como en RQ2, el
protocolo está \textbf{congelado}: se fijan de antemano los modelos, la
temperatura, el número de intentos, la versión del \emph{prompt}, el oráculo y
la composición del conjunto evaluado. El único componente intercambiable es el
\emph{proveedor} de respuestas, cuyo modo oficial es \lstinline|live| (la
invocación real a la API); esa indirección permite, por ejemplo, sustituir las
llamadas por respuestas grabadas en las pruebas sin alterar el resto del
protocolo.

\section{Reproducibilidad}
La reproducibilidad del estudio no se enuncia como aspiración, sino que se apoya
en elementos concretos y verificables:
\begin{itemize}
  \item exigencia explícita de \lstinline|OPENAI_API_KEY| en modo \lstinline|live|;
  \item versionado del \emph{prompt} propagado a cada artefacto;
  \item versiones fijadas en \texttt{pom.xml} (Java 17, JavaParser 3.26.4,
  JUnit 5.11.4, Gson 2.11.0);
  \item artefactos persistentes auditables sin ejecutar el código.
\end{itemize}
